\subsection{A two-mark scalar and its separated color sum}\label{c4:two-mark}

The diameter tails in \eqref{c4:Tail} determine the first length moment.
They do not control polygons that accumulate many edges in a small region.
For that purpose we will prove
\[
 \sum_{\operatorname{diam}\gamma\le r}|\gamma|^2x_*^{|\gamma|}
       \le r^{2/3+o(1)}.
\]
The sum is over the same unoriented translation classes as in
Section~\ref{c4:tails}.  We obtain it by marking two edges, bounding the
weight through a prescribed pair, and summing over their displacement.

There are three new ingredients.  First, the two-mark observable has a
Pfaffian denominator whose growth must be retained when comparing angular
and radial expressions.  Second, its ten screening charges require a
coercivity estimate on a doubled charge space.  Third, unequal arc lengths
leave an intermediate radial region in which only one color sector is
suppressed.  The partial stationary trace in
Section~\ref{c4:partial-trace} controls that region.  We begin with the
finite scalar identity, before making any asymptotic or positivity claim.


Again \(\lambda=d_*,\ Q=e^{\pi i/8},\ q=Q^2,\ \eta=-i\), with the normalized vacuum \(\Psi=\Psi_\eta\) of construction \(p=1\). Take the cyclic site list written in order \((P,L,T,R)\), where \(P,T\) are single distinguished sites and \(L,R\) the intervening blocks (seam immediately before \(P\)). Here \(P,T\) denote marks, not momentum, period or transfer symbols. We use the scalar
\[
 \mathscr H=
 \sum_{s=(s_L,s_R)} \eta^{S_L}\,
       \Psi(z)_{+\,s_L\,-\,s_R}\ \Psi(z^{-1})_{-\,s_L\,+\,s_R},
       \qquad S_L=\sum_{j\in L}s_j .
\]
In this calculation write \(x_j=z_j^2\), and use even total size \(N=2+|R|+|L|\). Put
\(\widetilde D(X)=D_1(X)\prod z_j^{-(N-1)}\), and \(d_0=Q^3+Q^{-3}=\sqrt{2-\sqrt2}\).
For complex log-site arguments use consistent square roots and set
\(P_l(\tau)=2\sin(\pi(l+\tau)/8)\), where \(q^\tau=x/y\) with the chosen log difference.

Color \(R,L\) separately by \(-,0,+\) (values \(-1,0,1\)), with total color sum equal in the two sectors. Give colors \(\pm\) local weight \(d_0\), color 0 weight \(1\), and each mark weight \(d_0\). Multiply a pair at \(x,y\) in any fixed order by
\[
          (-1)^\sigma\prod_{l=0}^7 P_l(\tau)^{\ell_l},
\]
as follows. Use \(e_l\) for the unit coordinate at \(l\bmod 8\). For two nonzero colors \(s,t\) on the same block use
\[
 \begin{cases}
 e_2+e_3+e_{-2}+e_{-3},&s=t,\\
 e_{-2}+e_{-3}-e_0-e_1,&s=-,\ t=+,\\
 e_2+e_3-e_0-e_{-1},&s=+,\ t=-.
 \end{cases}
\]
Across blocks for two nonzero colors use \(0\) when equal and \(2(e_{2s}+e_{3s})\) when opposite. For a mark first and a nonzero color \(t\) second use \(e_{-2t}+e_{-3t}\) (order reversal reflects indices); two marks use 0. To complete the exponents at a zero-colored site \(x\), replace it by the ordered pair consisting of a \(-\) site at \(q^{-1}x\) and a \(+\) site at \(qx\). Both split sites are present, so their shifted exponent vectors add; the denominator contributes the correction \(-2e_4\). Thus
\[
 \ell(0,t)=\mathcal T^{-1}\ell(-,t)+\mathcal T\ell(+,t)-2e_4,\qquad
 \ell(s,0)=\mathcal T\ell(s,-)+\mathcal T^{-1}\ell(s,+)-2e_4,\qquad \mathcal T e_l=e_{l+1}.
\]
Signs have \(\sigma=1\) just on within-block pairs directed \(-\to0,\ 0\to+,\ +\to-\), and on mark-color pairs with nonzero color. Denote the resulting sum of products by \(\mathscr N\). Then
\begin{equation}
                          \mathscr H=\mathscr N/\widetilde D(X)^2.                  \label{c4:Pair-sum}
\end{equation}
The identity is an identity of rational functions in the original
site-root variables.  At collisions we continue the complete quotient;
individual color terms may have poles.  To prove the identity, we first
identify two vacuum components, then resolve the block operators into
their one-dimensional eigenspaces, and finally remove vacant colors by
fusion.

\paragraph{Vacuum fusion and saturated components.}

First note two polynomial consequences of the vacuum construction. One is the following \(D_1\) fusion, valid in either parity:
\begin{equation}
      D_1(Y,q^{-1}x,qx)=\sqrt2 x\prod_{t\in Y}(t+x)\ D_1(Y,x).
                                                                    \label{c4:D1-fusion}
\end{equation}
Indeed by symmetry and the \(q^3\)-recursion (place the pair first), at two spectators \(u,q^3u\) both sides reduce with matching factors since
\[
 \frac{\prod_{\pm}(q^{\pm1}x-q^{-3}u)(q^{\pm1}x-q^{-2}u)}
             {(x-q^{-3}u)(x-q^{-2}u)}=(x+q^3u)(x+u).
\]
Induct in \(m=|Y|\), starting with \(D_1(a,b)=a+b\). Degree in a spectator is at most \(m+1\); zero and leading infinity match by deletion, and the \(2(m-1)\) pair nodes just found suffice when \(m\ge2\). For \(m=1\) test additionally at that spectator \(q^2x\): deletion of \(q^{-1}x,q^2x\) on the left gives \(2(1+q)(1+q^3)x^3=\sqrt2(1+q^2)^2x^3\).

The other consequence for size \(2k\) is a saturated component (\(s=\pm1\)):
\[
 \Psi(z)_{s^k\,(-s)^k}
     =(d_0 q^s)^k \ \frac{\prod_{i<j,\ {\rm same\ spin}}S_2(x_i,x_j)S_3(x_i,x_j)}
                                    {\widetilde D(X)},\quad
 S_l(x,y)=\frac{z_y}{z_x}Q^{-l}-\frac{z_x}{z_y}Q^l=-iP_l(\tau).
\]
Indeed stripped \(V_{2k}\) has degree at most \(2k-2\) per occupied variable; same-spin swaps act by scalar \(w\), hence zeros at \(x_j=q^2x_i,q^3x_i\) for ordered pairs within either run (bring them adjacent within the run). The resulting product saturates the bound, giving the formula up to its stated constant. At \(x_k=x,\ z_{k+1}=Q^3 z_k\) adjacent deletion gives \(q^s\) times the smaller normalized vacuum component. Here \(\widetilde D(X)/\widetilde D(X_{\rm red})=d_0\prod_{t\in X_{\rm red}}S_{-2}(x,t)S_{-3}(x,t)\) directly by the \(D_1\) recursion (a sign per spectator, even number); the extra pair factors upstairs are precisely the displayed spectator product using \(S_{l+8}=-S_l\). This proves the constant by induction.

\paragraph{Resolving the two blocks.} On each block alone use the open-row operators \(M_{p',p'}\) with \(p'=-1\) on \(R\), \(p'=+1\) on \(L\), and generic auxiliary arguments. In block spin basis they are triangular: at the first change between input and output the output-minus-input spin equals the new intermediate auxiliary spin minus \(p'\), by charge conservation. Thus lex order can only increase for \(p'=-\), decrease for \(p'=+\). Diagonal eigenvalues assigned to a color list are products of \(w,v,W\), respectively for colors \(-,0,+\) meaning actual diagonal-list spins \(p',0,-p'\). These spectra are simple generically: local numerator zero exponents in the auxiliary square relative to the site square are \(\{-2,-3\},\{0,-3\},\{0,-1\}\), disjoint from other generic site orbits. Insert the identity in the form of the tensor product spectral resolution in \(\mathscr H\) between the vacuum components; the block charges commute and require precisely the color-sum constraint. The factor \(\eta^{S_L}\) is scalar at those charges.

Each term with a zero color \(x\) equals the corresponding term on a larger system replacing that site by \(q^{-1}x,qx\) in order (roots \(Q^{-1}z,Qz\)), colored \(-,+\). Indeed positive-order pair fusion inserts the map \(A\) on the lower vacuum, and \(A^t\) on the inverse-site vacuum gives the shorter inverse vacuum, by the same vacuum argument as in the marked-row reductions. The spatial transport fuses the open row on this subspace, and all block charges agree. The chosen spectral label lifts uniquely: zeros for the two local numerators now use one choice from \(\{\{-3,-4\},\{-1,-4\},\{-1,-2\}\}\) and one from \(\{\{-1,-2\},\{1,-2\},\{1,0\}\}\), with common denominator exponents \(1,2,3,4\); all nine multisets are distinct modulo 8. The required fused vacancy zeros thus match just \((-,+)\). So the spectral projections intertwine on the inserted subspace. These arguments apply also to several simultaneously split generic sites, by the same simple spectra, rational specialization, and nonvanishing of denominators generically there by \eqref{c4:D1-fusion}. This reduces all projection terms to full occupancy.

\paragraph{Evaluating a fully occupied term.} Now work first on a generic fully occupied term (then specialize as above). Sort the variables on each block separately to put \(-\) before \(+\), retaining the two internal orders. Simultaneous covariance is by the ordinary exchanges, inverse-transposed in the inverse-site vacuum. In the sorted order the left eigenvector in each block is just its indicated spin bra: it is the lexicographic extreme of its charge with no inputs other than itself to that row. In reversed-run order (\(+\) before \(-\)) the right eigenvector similarly is the bare ket. Transport it to sorted order by the block-exchange braid; contracting with the inverse-site vacuum now reads the latter in reversed-run order, by transposed exchanges. The normalizing left-right pairing on each block is
\(\prod_{a\ {\rm of\ color}\ -,\, b\ {\rm of\ color}\ +} W(z_b/z_a)\).
Indeed move the minus-colored sites from right to left one at a time over the plus run. Each moved spin has to be \(p'\) on arrival (it stays fixed from then on); total charge forces the plus run still saturated then, leaving only the spin-preserving-along-labels swap at every crossing, of entry \(W\).

Thus use the two saturated amplitudes above: the lower one after starting its circle at the sorted \(R_+\) run (to read \(+\) first), the inverse upper one at \(R_-\) (to read \(-\) first), or at \(P\) if the indicated run is empty. Recovering the original amplitudes gives a factor \(\eta^{|R_+|-|R_-|}\) canceling \(\eta^{S_L}\). The constants now give \(d_0\) per site, and the denominators are \(\widetilde D(X)\widetilde D(X^{-1})\). For like colors on the same block the product contributes the \(S_l\) at \(l=\pm2,\pm3\); for unlike colors in order minus then plus the normalization gives \(1/W=S_{-2}S_{-3}/(S_0 S_1)\). For a pair first on \(R\) then \(L\), equal actual spins occur at opposite colors \(s,-s\); the two saturated orders place that plus-spin pair first as \(R,L\) in lower and \(L,R\) in inverse, or oppositely for minus spins. Thus the factors are \(S_{2s}^2 S_{3s}^2\); similarly with arguments reversed. Putting either mark first versus a color \(t\), it meets that color with like spin on exactly one side in order giving \(S_{-2t} S_{-3t}\) (two reversed-argument signs cancel). This proves the nonzero-color rules after \(S_l=-iP_l,\ P_{l+8}=-P_l\), with even-size inversion of \(\widetilde D\) as in the Pfaffian formula below.

\paragraph{Restoring the vacant colors.} Finally \eqref{c4:D1-fusion}, normalized, gives \(\sqrt2\prod_y P_4(\tau_{xy})\) per split in each denominator. The internal new pair contributes \(1/W(Q^2)=2+\sqrt2\), so local color-zero weight after removing the constants is \(d_0^2(2+\sqrt2)/2=1\). All other factors just shift their lags as in the rules and pay \(2e_4\) from the denominator. For clarity, when expanding the first of a pair the new parity adds the two previous parities and \(\ell_0(-,t)+\ell_7(+,t)\) from wrapping the index range; for expanding the second it uses \(\ell_7(s,-)+\ell_0(s,+)\) instead. Applying this to the displayed base tables gives the within-block sign rows \((0,1,0),(0,0,1),(1,0,0)\), zero parities across blocks, and positive sign at a mark paired with a vacancy. This completes \eqref{c4:Pair-sum}.
